\begin{pastexam}[label={px:ntu-113-bme-ds-01}]{臺大生機所 113 學年度 第 1 題}
（問答）

\textbf{(10\%)} For a \textbf{linked list} implementation of a \textbf{stack}, what do you think the
easiest and most efficient implementation of the \textbf{push} (to the front of the linked list or to
the end) and \textbf{pop} (from the front of the linked list or from the end) operations? Please
justify your answer.
\end{pastexam}

\begin{pastsoln}
答：push 和 pop 都在串列的\textbf{前端}做，兩個都是 $\Theta(1)$，而且只用單向串列加一個
$\mathit{head}$ 指標就夠了（資結書第~6.1~節、虛擬碼~6.2）。

\textbf{前端的做法。} 設串列是單向的，$L.\mathit{head}$ 指向第一個節點。
\begin{itemize}
\item push$(x)$：$x.\mathit{next} \gets L.\mathit{head}$，$L.\mathit{head} \gets x$。 改兩個指標。
\item pop：$x \gets L.\mathit{head}$，$L.\mathit{head} \gets x.\mathit{next}$，回傳 $x$。 改一個指標
（串列空的時候回報 underflow）。
\end{itemize}
不管串列有多長，都只做固定幾步，$\Theta(1)$。 最後 push 進去的節點一定在最前面，pop 拿到的
正是它，符合後進先出。

\textbf{為什麼不放在尾端。} 堆疊的 push 和 pop 必須在同一端，否則就變成佇列了。 如果兩個都在
尾端：
\begin{itemize}
\item push 到尾端：要先走到最後一個節點，$\Theta(n)$；多存一個 $\mathit{tail}$ 指標可以降到
$\Theta(1)$。
\item pop 尾端：拿掉最後一個節點後，要把\textbf{倒數第二個}節點的 $\mathit{next}$ 設成 \NIL，並把
$\mathit{tail}$ 移到它。 單向串列的節點不知道前一個是誰，只能從頭走過去找，$\Theta(n)$；
就算有 $\mathit{tail}$ 也一樣。
\end{itemize}
要讓尾端的 pop 也是 $\Theta(1)$，得改用雙向串列（每個節點多一個 $\mathit{prev}$），再維護
$\mathit{tail}$，空間和程式都比較多。 所以前端最簡單、也最有效率：兩個操作都是 $\Theta(1)$，
不需要額外的指標。
\end{pastsoln}
